\documentclass[a4paper,landscape,8pt]{extarticle}
\usepackage[landscape,margin=0.7cm,top=0.5cm,bottom=0.7cm]{geometry}
\usepackage[utf8]{inputenc}
\usepackage[T1]{fontenc}
\usepackage[ngerman]{babel}
\usepackage{helvet}
\renewcommand{\familydefault}{\sfdefault}
\usepackage{xcolor}
\usepackage{tikz}
\usetikzlibrary{positioning, shapes.geometric}
\usepackage{enumitem}
\usepackage{multicol}
\usepackage{array}
\usepackage{fancyhdr}
\usepackage{amssymb}

% Farben
\definecolor{privacyblue}{HTML}{5856D6}
\definecolor{datagreen}{HTML}{34C759}
\definecolor{permissionorange}{HTML}{FF9500}
\definecolor{trackingred}{HTML}{FF3B30}
\definecolor{lightgray}{HTML}{F5F5F5}
\definecolor{bordergray}{HTML}{BBBBBB}
\definecolor{warnred}{HTML}{C62828}
\definecolor{appleblue}{HTML}{007AFF}

\pagestyle{fancy}
\fancyhf{}
\renewcommand{\headrulewidth}{0pt}
\fancyfoot[C]{\footnotesize Seite \thepage{} von 3}

\setlength{\parindent}{0pt}
\setlength{\parskip}{2pt}

% Kompakte Listen
\setlist{nosep, leftmargin=1em, itemsep=0pt, topsep=1pt}

\begin{document}

% ============================================================================
% SEITE 1: INFORMATIONSBLATT
% ============================================================================

\noindent{\Large\bfseries Datenschutz \& Privatsphäre: Wer weiß was über dich?} \hfill
{\footnotesize\textbf{Lernziele:} Datensammlung verstehen | Berechtigungen prüfen | Privatsphäre schützen}

\vspace{1mm}
\hrule
\vspace{1mm}

% Drei-Spalten-Layout
\begin{minipage}[t]{0.32\linewidth}
\begin{center}
\colorbox{privacyblue!10}{%
\begin{minipage}{0.95\linewidth}
\vspace{2mm}
\begin{center}
{\Large\bfseries\textcolor{privacyblue}{Was sind \glqq Daten\grqq{}?}}\\[1pt]
{\footnotesize\itshape Alles, was über dich gespeichert wird}
\end{center}
\vspace{2mm}
\end{minipage}%
}
\end{center}

\vspace{1mm}

\textbf{Welche Daten gibt es?}
\begin{itemize}
\item \textbf{Name, Adresse, Geburtsdatum}
\item \textbf{E-Mail-Adresse, Telefonnummer}
\item \textbf{Standort} (wo du bist)
\item \textbf{Fotos und Videos}
\item \textbf{Kontakte} (dein Adressbuch)
\item \textbf{Surfverhalten} (welche Webseiten)
\item \textbf{Einkäufe} (was du kaufst)
\item \textbf{Gesundheitsdaten} (Schritte, Puls)
\end{itemize}

\textbf{Analogie:} Daten sind wie \textbf{Puzzleteile} -- einzeln harmlos, zusammen ergeben sie ein komplettes Bild von dir.

\vspace{2mm}

\textbf{Wer sammelt Daten?}
\begin{itemize}
\item Apps auf deinem Handy
\item Webseiten, die du besuchst
\item Online-Shops
\item Soziale Netzwerke
\item Werbefirmen
\end{itemize}

\end{minipage}%
\hfill%
\begin{minipage}[t]{0.32\linewidth}
\begin{center}
\colorbox{datagreen!10}{%
\begin{minipage}{0.95\linewidth}
\vspace{2mm}
\begin{center}
{\Large\bfseries\textcolor{datagreen}{Warum wird gesammelt?}}\\[1pt]
{\footnotesize\itshape Nicht immer böse, aber oft unnötig}
\end{center}
\vspace{2mm}
\end{minipage}%
}
\end{center}

\vspace{1mm}

\textbf{Gute Gründe:}
\begin{itemize}
\item Navigation braucht deinen \textbf{Standort}
\item Fotos-App braucht Zugriff auf \textbf{Kamera}
\item WhatsApp braucht deine \textbf{Kontakte}
\end{itemize}

\textbf{Fragwürdige Gründe:}
\begin{itemize}
\item Werbung auf dich \textbf{zuschneiden}
\item Dein \textbf{Verhalten analysieren}
\item Daten an andere \textbf{verkaufen}
\end{itemize}

\textbf{Analogie:} Wie ein Geschäft, das \textbf{heimlich notiert}, was du dir anschaust -- und diese Info verkauft.

\vspace{2mm}

\textbf{Das Geschäftsmodell:}\\
Viele Apps sind \glqq kostenlos\grqq{} -- aber du bezahlst mit deinen \textbf{Daten}.

\begin{center}
\fcolorbox{datagreen}{datagreen!8}{%
\begin{minipage}{0.9\linewidth}
\vspace{2mm}
\textbf{Faustregel:}\\[2pt]
Wenn etwas kostenlos ist, bist oft \textbf{du das Produkt}.
\vspace{2mm}
\end{minipage}}
\end{center}

\end{minipage}%
\hfill%
\begin{minipage}[t]{0.32\linewidth}
\begin{center}
\colorbox{warnred!10}{%
\begin{minipage}{0.95\linewidth}
\vspace{2mm}
\begin{center}
{\Large\bfseries\textcolor{warnred}{Was kann schiefgehen?}}\\[1pt]
{\footnotesize\itshape Risiken bei zu viel Datenfreigabe}
\end{center}
\vspace{2mm}
\end{minipage}%
}
\end{center}

\vspace{1mm}

\textbf{Mögliche Probleme:}
\begin{itemize}
\item \textbf{Nervige Werbung} -- überall verfolgt dich die gleiche Anzeige
\item \textbf{Preismanipulation} -- teurere Angebote für dich
\item \textbf{Datenlecks} -- deine Daten werden gestohlen
\item \textbf{Identitätsdiebstahl} -- jemand gibt sich als du aus
\item \textbf{Einbrecher} -- wissen, wann du nicht zu Hause bist
\end{itemize}

\vspace{2mm}

\textbf{Beispiel Standort:}\\
Wenn du ständig deinen Standort teilst, wissen Fremde:
\begin{itemize}
\item Wo du wohnst
\item Wo du arbeitest
\item Wann du nicht zu Hause bist
\item Welche Wege du gehst
\end{itemize}

\vspace{2mm}

\begin{center}
\fcolorbox{warnred}{warnred!8}{%
\begin{minipage}{0.9\linewidth}
\vspace{2mm}
\textbf{Grundregel:}\\[2pt]
Gib nur die Daten frei, die eine App \textbf{wirklich braucht}.
\vspace{2mm}
\end{minipage}}
\end{center}

\end{minipage}

\vspace{1mm}

% Zusammenfassung
\begin{center}
\fcolorbox{bordergray}{lightgray}{%
\begin{minipage}{0.98\linewidth}
\centering
\vspace{1.5mm}
\textbf{Merke:} Daten sind wertvoll. Prüfe bei jeder App: \textbf{Braucht sie das wirklich?} Im Zweifel: \textbf{Nein} sagen.
\vspace{1.5mm}
\end{minipage}%
}
\end{center}

\newpage

% ============================================================================
% SEITE 2: EINSTELLUNGEN & BERECHTIGUNGEN
% ============================================================================

\begin{center}
{\LARGE\bfseries Berechtigungen \& Datenschutz-Einstellungen}\\[3pt]
{\normalsize So kontrollierst du, wer was darf}
\end{center}

\vspace{1mm}
\hrule
\vspace{2mm}

\begin{multicols}{2}

\section*{\textcolor{permissionorange}{App-Berechtigungen}}

\textbf{Was sind Berechtigungen?} Die \textbf{Erlaubnis}, die eine App braucht, um bestimmte Dinge zu tun.

\textbf{Analogie:} Wie ein \textbf{Schlüssel} -- du entscheidest, wer welchen Schlüssel bekommt.

\vspace{2mm}

\textbf{Wichtige Berechtigungen:}

\renewcommand{\arraystretch}{1.3}
\begin{tabular}{@{}p{2.2cm}p{5.3cm}@{}}
\textbf{Standort} & Weiß, wo du bist \\
\textbf{Kamera} & Kann Fotos/Videos machen \\
\textbf{Mikrofon} & Kann zuhören \\
\textbf{Fotos} & Zugriff auf deine Bilder \\
\textbf{Kontakte} & Sieht dein Adressbuch \\
\textbf{Kalender} & Sieht deine Termine \\
\textbf{Gesundheit} & Zugriff auf Gesundheitsdaten \\
\end{tabular}
\renewcommand{\arraystretch}{1}

\vspace{2mm}

\textbf{Berechtigungen prüfen:}\\
Einstellungen → \textbf{Datenschutz \& Sicherheit}

Dort siehst du für jede Berechtigung, \textbf{welche Apps} sie haben.

\vspace{2mm}

\textbf{Standort-Optionen:}
\begin{itemize}
\item \textbf{Nie} -- App bekommt keinen Standort
\item \textbf{Beim Verwenden} -- Nur wenn App offen ist
\item \textbf{Immer} -- Auch im Hintergrund (selten nötig!)
\end{itemize}

\vspace{2mm}

\begin{center}
\fcolorbox{permissionorange}{permissionorange!8}{%
\begin{minipage}{0.9\linewidth}
\vspace{2mm}
\textbf{Tipp:}\\[2pt]
Wähle bei Standort am besten \textbf{\glqq Beim Verwenden\grqq{}} -- nicht \glqq Immer\grqq{}.
\vspace{2mm}
\end{minipage}}
\end{center}

\columnbreak

\section*{\textcolor{trackingred}{Tracking verhindern}}

\textbf{Was ist Tracking?} Apps und Webseiten \textbf{verfolgen dich} über verschiedene Seiten hinweg.

\textbf{Analogie:} Wie ein \textbf{Detektiv}, der dir überall folgt und notiert, was du machst.

\vspace{2mm}

\textbf{App-Tracking ablehnen:}
\begin{enumerate}
\item Einstellungen → Datenschutz \& Sicherheit
\item \textbf{Tracking} antippen
\item \textbf{\glqq Apps erlauben, Tracking anzufordern\grqq{}} \textbf{ausschalten}
\end{enumerate}

$\rightarrow$ Apps können dich nicht mehr über andere Apps verfolgen!

\vspace{2mm}

\textbf{Safari-Datenschutz:}
\begin{enumerate}
\item Einstellungen → Safari
\item \textbf{\glqq Cross-Sitetracking verhindern\grqq{}} einschalten
\item \textbf{\glqq Alle Cookies blockieren\grqq{}} optional
\end{enumerate}

\vspace{3mm}

\section*{\textcolor{appleblue}{Apples Datenschutz-Funktionen}}

\textbf{Apple betont Datenschutz:}
\begin{itemize}
\item \textbf{App-Datenschutzbericht} -- zeigt, was Apps tun
\item \textbf{E-Mail-Datenschutz} -- versteckt deine IP
\item \textbf{Privates Surfen} -- Safari speichert nichts
\item \textbf{Anmelden mit Apple} -- keine echte E-Mail verraten
\end{itemize}

\vspace{2mm}

\begin{center}
\fcolorbox{appleblue}{appleblue!8}{%
\begin{minipage}{0.9\linewidth}
\vspace{2mm}
\textbf{Tipp: \glqq Mit Apple anmelden\grqq{}}\\[2pt]
Wenn eine App das anbietet: Nutze es! Apple erstellt eine \textbf{Wegwerf-E-Mail}, die zu dir weiterleitet.
\vspace{2mm}
\end{minipage}}
\end{center}

\end{multicols}

\newpage

% ============================================================================
% SEITE 3: ÜBUNGEN & CHECKLISTE
% ============================================================================

\begin{center}
{\LARGE\bfseries Übungen \& Checkliste}\\[3pt]
{\normalsize Prüfe deine Datenschutz-Einstellungen}
\end{center}

\vspace{1mm}
\hrule
\vspace{2mm}

\begin{multicols}{2}

\textbf{\large Teil A: Braucht die App das wirklich?}

Sinnvoll ($\checkmark$) oder fragwürdig (?)?

\vspace{2mm}

\renewcommand{\arraystretch}{1.4}
\begin{tabular}{@{}p{5cm}p{2cm}cc@{}}
App & Berechtigung & $\checkmark$ & ? \\
\hline
Karten-App & Standort & $\square$ & $\square$ \\
Taschenlampe & Kontakte & $\square$ & $\square$ \\
WhatsApp & Kontakte & $\square$ & $\square$ \\
Spiele-App & Mikrofon & $\square$ & $\square$ \\
Kamera-App & Kamera & $\square$ & $\square$ \\
Wetter-App & Standort & $\square$ & $\square$ \\
Rezepte-App & Standort & $\square$ & $\square$ \\
\end{tabular}
\renewcommand{\arraystretch}{1}

\textit{\footnotesize Lösungen: $\checkmark$, ?, $\checkmark$, ?, $\checkmark$, $\checkmark$, ?}

\vspace{4mm}

\textbf{\large Teil B: Richtig oder Falsch?}

\vspace{2mm}

\renewcommand{\arraystretch}{1.4}
\begin{tabular}{@{}p{7.5cm}cc@{}}
& R & F \\
1. Kostenlose Apps sammeln oft Daten. & $\square$ & $\square$ \\
2. \glqq Standort: Immer\grqq{} ist die sicherste Option. & $\square$ & $\square$ \\
3. Man kann App-Tracking komplett ausschalten. & $\square$ & $\square$ \\
4. \glqq Mit Apple anmelden\grqq{} schützt deine E-Mail. & $\square$ & $\square$ \\
5. Berechtigungen kann man nachträglich ändern. & $\square$ & $\square$ \\
\end{tabular}
\renewcommand{\arraystretch}{1}

\textit{\footnotesize Lösungen: R, F, R, R, R}

\vspace{4mm}

\textbf{\large Teil C: Praktische Übungen}

Führe diese Aktionen auf deinem iPhone aus:

\vspace{2mm}

\begin{itemize}[label=$\square$, itemsep=3pt]
\item Öffne \textbf{Einstellungen → Datenschutz \& Sicherheit}
\item Tippe auf \textbf{Ortungsdienste} -- welche Apps haben Zugriff?
\item Tippe auf \textbf{Tracking} -- ist es ausgeschaltet?
\item Prüfe, welche Apps auf \textbf{Mikrofon} zugreifen dürfen
\item Prüfe, welche Apps auf \textbf{Kamera} zugreifen dürfen
\item Öffne \textbf{Einstellungen → Safari} und prüfe Datenschutz-Optionen
\end{itemize}

\columnbreak

\textbf{\large Teil D: Aufräumen}

Gibt es Apps mit Berechtigungen, die sie nicht brauchen?

\vspace{2mm}

\begin{tabular}{@{}p{4cm}p{4cm}@{}}
App: & Berechtigung entfernt: \\
\dotfill & \dotfill \\[3mm]
\dotfill & \dotfill \\[3mm]
\dotfill & \dotfill \\
\end{tabular}

\vspace{4mm}

\textbf{\large Checkliste: Das habe ich geprüft}

\vspace{2mm}

\begin{center}
\fcolorbox{privacyblue}{privacyblue!10}{%
\begin{minipage}{0.95\linewidth}
\vspace{2mm}
\textbf{\textcolor{privacyblue}{Grundwissen}}
\begin{itemize}[label=$\square$, itemsep=2pt]
\item Ich verstehe, warum \textbf{Datenschutz wichtig} ist
\item Ich weiß, welche \textbf{Daten} gesammelt werden können
\item Ich weiß: Kostenlos = oft mit Daten bezahlt
\end{itemize}
\vspace{2mm}
\end{minipage}}
\end{center}

\vspace{2mm}

\begin{center}
\fcolorbox{permissionorange}{permissionorange!10}{%
\begin{minipage}{0.95\linewidth}
\vspace{2mm}
\textbf{\textcolor{permissionorange}{Berechtigungen}}
\begin{itemize}[label=$\square$, itemsep=2pt]
\item Ich habe \textbf{Ortungsdienste} geprüft
\item Ich habe \textbf{Kamera/Mikrofon}-Zugriffe geprüft
\item Ich habe unnötige Berechtigungen \textbf{entzogen}
\end{itemize}
\vspace{2mm}
\end{minipage}}
\end{center}

\vspace{2mm}

\begin{center}
\fcolorbox{trackingred}{trackingred!10}{%
\begin{minipage}{0.95\linewidth}
\vspace{2mm}
\textbf{\textcolor{trackingred}{Tracking}}
\begin{itemize}[label=$\square$, itemsep=2pt]
\item \textbf{App-Tracking} ist ausgeschaltet
\item \textbf{Cross-Sitetracking} in Safari ist verhindert
\end{itemize}
\vspace{2mm}
\end{minipage}}
\end{center}

\vspace{2mm}

\begin{center}
\fcolorbox{bordergray}{white}{%
\begin{minipage}{0.95\linewidth}
\vspace{2mm}
\textbf{Meine Notizen}\\[2pt]
Apps, denen ich Berechtigungen entzogen habe:\\[2mm]
\dotfill
\vspace{2mm}
\end{minipage}}
\end{center}

\end{multicols}

\vfill

\begin{center}
\footnotesize\textit{Stand: \today{} -- Prüfe regelmäßig deine Berechtigungen. Weniger ist mehr!}
\end{center}

\end{document}
